% v2_lemma_v3.tex -- round 374, lemma-first attack on V2.
% Outcome: REDUCED to the strictly smaller V3' (15-node
% Pruefer-step profile).  Builds with pdflatex.
% Companion machine check: verify_v2_lemma.py.
% NO RH CLAIM.  Research documentation, not a theorem of RH.
\documentclass[11pt]{article}

\usepackage[a4paper,margin=2.7cm]{geometry}
\usepackage{amsmath,amssymb,amsthm}
\usepackage{booktabs}
\usepackage{microtype}
\usepackage{xcolor}
\usepackage[colorlinks=true,linkcolor=blue!50!black,urlcolor=blue!50!black]{hyperref}

\newtheorem{lemma}{Lemma}
\newtheorem{prop}{Proposition}
\newtheorem{definition}{Definition}
\theoremstyle{remark}
\newtheorem{remark}{Remark}

\newcommand{\N}{\mathbb{N}}
\newcommand{\Q}{\mathbb{Q}}
\newcommand{\R}{\mathbb{R}}

\title{\bfseries The $V_2$-lemma, reduced\\[2mm]
\large Exact Pr\"ufer step formula, the contrast obstruction,
and the strictly smaller profile lemma $V_3'$}
\author{Stefan Hamann}
\date{August 28, 2026}

\begin{document}
\maketitle

\begin{abstract}\noindent
Round 374 attacks the named remainder $V_2$ of~\cite{xn,v2}
by the r372 Pr\"ufer-to-run dictionary (a theorem: sign-runs of
$v_2$ equal $\lfloor\theta/\pi\rfloor$-dwells of the
degree-accumulated Jacobi--Pr\"ufer phase).  $V_2$ itself is
neither proved nor refuted in the construction.  What is proved:
an exact two-point Wronskian formula for the discrete phase step
of the scaled pair $(v_{N-2},v_{N-3})$; that at the $+x$
$20\%$ mask the colouring is $\mathrm{sign}(w)\,\mathrm{sign}(v_2)$
($x$ does not flip); a finite contrast obstruction on
$\R^{24}$ (constant steps never realise the violator occupancy
for any remainder; a slow-then-fast block $(0.4^{21},2.8^{3})$
does, while $(1.0^{21},2.0^{3})$ does not);
and that a coherent $\times 8$ rescaling of the last twelve
Jacobi $\gamma$'s on $\chi_3$ at $w=9$ \emph{does} produce a
$v_2$-violator, so $V_2$ is not a theorem of abstract
Jacobi-positive chains.  A gap-then-spike weight adversary on
the real fold-grid at $8\times$ and $64\times$ never produces a
$V_2$-violator and always breaks Jacobi positivity.  The hydra-legal
reduction is the finite profile lemma $V_3'$: the last fourteen
Pr\"ufer steps at the mask lie in the region $\mathcal{A}_{15}$
of sequences that are $V_2$-regular for \emph{every} incoming
remainder $\eta\in[0,\pi)$.  That region is strictly smaller
than $V_2$ (finite window, Wronskian/mesh, no product history).
On the pin windows the source steps lie in $\mathcal{A}_{15}$
(Path~A: $0$ remainder-violators; every $(1,1,1)$ that occurs
is $V_2$-regular).  The chain
$T_2'\Leftarrow g_i\ge 3/8\Leftarrow\mathrm{MED\text{-}CAP}
\Leftarrow X_n\Leftarrow V_2\Leftarrow V_3'$
stands, with $V_3'$ the remaining named lemma.  No RH claim.
\end{abstract}

\medskip\noindent
\textbf{Status.}\enspace
Lemmas~\ref{lem:wronskian}--\ref{lem:contrast} are proved
(algebraic identity; finite step-region dichotomy).
Lemma~\ref{lem:xconst} is proved (mask geometry).
Lemma~\ref{lem:not-abstract} is a machine-checked negative
(an $8\times$ last-$12$ $\gamma$-scale on $\chi_3$ $w=9$ is a
$v_2$-violator).  Lemma~\ref{lem:weight-adv} is a machine-checked
negative (weight-cluster on FRAME-A $w=9$ never violates, and
breaks positivity).  Proposition~\ref{prop:v3} names $V_3'$ and
records Path~A on the pin windows.  The hydra test of~\cite{v2}
is met: $V_3'$ is strictly more elementary than $V_2$.
No RH claim.

\section{Recollection}\label{sec:rec}

$V_2$ of~\cite{xn}~\S5: if the last three $x$-order run lengths
of $\mathrm{sign}(v_2\,w\,x)$ at the $20\%$ $x$-mask are
$(1,1,1)$, then among the four preceding at least two are
$\le 2$.  The five-run pattern $(\ldots,r,r,1,1,1)$ with
$r\ge 3$ is a strictly stronger combinatorial statement (the
spike $(2^{8},3,3,1,1,1)$ is $V_2$-regular; gate G2) and is
not the named remainder.  The r365 anatomy~\cite{v2}: $v_2$ is the
degree-$(N-2)$ arithmetic-comb chain polynomial on the
smooth-border nodes (source Jacobi, not mesh-only); $w$ is the
explicit smooth-fold arm; a $2$-regular factor does not imply
$V_2$; the $134$-window surface has $0$ violators.  The r372
dictionary is a theorem: $\lfloor\theta/\pi\rfloor$ runs of the
chain-accumulated Jacobi--Pr\"ufer phase equal the sign-runs of
$v_2$ (181/181, XOR 181/181, interior progress $\ge 2.062$).
Run lengths \emph{are} phase dwell times.  This note attacks
$V_2$ with that dictionary, under the lemma-first charter
(exactly one named lemma; exit proved / refuted / reduced).

\section{The exact discrete step}\label{sec:step}

Write $v_n$ for the scaled monic chain value of degree $n=N-2$
and $v_{n-1}$ for the previous scaled value, as produced by the
three-term recurrence of~\cite{xn}~\S2 (coefficients
$(\hat\alpha_k,\gamma_{k+1})$, log-scale $L_s$).  Let
$\varphi(x)=\mathrm{atan2}(v_n(x),v_{n-1}(x))$ and
$r(x)=\sqrt{v_n(x)^2+v_{n-1}(x)^2}$.

\begin{lemma}[Wronskian step]\label{lem:wronskian}
For any two points $x,y$ at which $r(x)r(y)>0$,
\[
\sin\bigl(\varphi(y)-\varphi(x)\bigr)
 \;=\;
 \frac{v_n(y)\,v_{n-1}(x)-v_{n-1}(y)\,v_n(x)}{r(x)\,r(y)}.
\]
Equivalently, writing $\mathbf{u}=(v_n,v_{n-1})$, the identity
\[
\bigl(\mathbf{u}(x)\wedge\mathbf{u}(y)\bigr)^2
 +
 \bigl(\mathbf{u}(x)\cdot\mathbf{u}(y)\bigr)^2
 \;=\;
 \lvert\mathbf{u}(x)\rvert^2\,\lvert\mathbf{u}(y)\rvert^2
\]
holds in any field of characteristic not $2$ (Lagrange).
\end{lemma}

\begin{proof}
The two-dimensional Lagrange identity
$(a d-b c)^2+(a c+b d)^2=(a^2+b^2)(c^2+d^2)$.  The sine formula
is the imaginary part of
$e^{i(\varphi(y)-\varphi(x))}=
(\mathbf{u}(y)\cdot\overline{\mathbf{u}(x)})/\lvert\mathbf{u}(x)\rvert\lvert\mathbf{u}(y)\rvert$
in the identification $\mathbf{u}\leftrightarrow v_n+i v_{n-1}$.
\end{proof}

The degree-accumulated Pr\"ufer phase $\theta$ of the r372
dictionary is the lift of $\varphi$ along the recurrence in the
degree direction, then unwrapped along $x$.  Its spatial steps
$\Delta\theta_i=\theta(x_{i+1})-\theta(x_i)$, reduced to
$(-\pi,\pi]$ and oriented so that $\sum\Delta\theta_i>0$, are
the dwell increments: a $\pi$-bin change (a $v_2$ sign-flip)
occurs between consecutive nodes iff those nodes occupy
distinct $\lfloor\theta/\pi\rfloor$ levels.  Lemma~\ref{lem:wronskian}
identifies the sine of the \emph{principal} step with the
scaled Wronskian; the accumulated $\theta$-step agrees with the
principal step up to the $2\pi$-unwrap, which the companion
script gates on the pin windows.

\begin{lemma}[$+x$ mask, $x$ constant]\label{lem:xconst}
On the last nine bulk-border nodes at the $+x$ $20\%$ hull mask,
$x>0$.  Consequently
$\mathrm{sign}(v_2 w x)=\mathrm{sign}(v_2)\,\mathrm{sign}(w)$
there (the r365 XOR identity, with no $x$-flip).
\end{lemma}

\begin{proof}
The mask keeps the central $60\%$ of the hull.  The $+x$ edge
is at $x\approx +0.60$ (hull near $[-1,1]$).  The last nine
bulk nodes sit in a mesh neighbourhood of that edge, on the
positive side.  The identity is the product rule.
\end{proof}

Thus $V_2$ at the load-bearing edge is a statement about the
XOR of two sign fields, of which $v_2$ is the Jacobi--Pr\"ufer
dwell sequence and $w$ is construction-explicit.  On the MAIN
family ($w$ a single bulk run) the XOR is vacuous and $V_2$ is
purely a $v_2$-dwell statement.

\section{The contrast obstruction}\label{sec:contrast}

A run of length $r$ is $r$ consecutive nodes in one $\pi$-bin,
hence $r-1$ steps totalling less than $\pi$.  For $r\ge 3$ the
mean interior step is $<\pi/r\le\pi/3\approx 1.047$.  Three
consecutive $1$-runs are three consecutive $\pi$-crossings: the
terminal steps must be large (empirically $\gtrsim 2.2$).  The
violator $(\ldots,r,r,1,1,1)$ is therefore a
\emph{slow-block then fast-block} in the step sequence, not a
uniform speed.

\begin{lemma}[contrast dichotomy]\label{lem:contrast}
Let $d=(\Delta_0,\ldots,\Delta_{23})\in(0,\pi)^{24}$ be a rigid
step sequence, and let $\mathrm{occupancy}(\eta,d)$ be the
$\lfloor\theta/\pi\rfloor$ run lengths of the walk
$\theta_0=\eta$, $\theta_{i+1}=\theta_i+\Delta_i$, on $25$
nodes.
\begin{enumerate}
\item Constant steps $\Delta_i\equiv\Delta$ for
  $\Delta\in\{0.8,1.0,1.2,\pi/2,1.8,2.0\}$: among $1001$
  remainders $\eta\in[0,\pi)$, $0$ $V_2$-violators and $0$
  five-run patterns $(r,r,1,1,1)$.
\item Slow-then-fast blocks of the form
  $(s^{21},f,f,f)$.  At $(s,f)=(0.4,2.8)$ there are $347/1001$
  official $V_2$-violators (and $674/1001$ five-run patterns).
  At $(s,f)=(1.0,2.0)$ both sets are empty.
\end{enumerate}
A rigid violator, when it occurs, is a slow block adjacent to
a fast block, not a constant speed.  The two gated profiles
separate at a step-ratio between $2$ and $7$.
\end{lemma}

The companion script re-proves the dichotomy by exact
occupancy counting (gate G3--G4).  It is a finite statement
about $\R^{24}$, independent of the source.  The Jacobi
coefficients $(\hat\alpha_k,\gamma_k)$ do \emph{not} depend on
$x$: there is no ``coefficient jump at the mask''.  The
$x$-variation of the step is entirely the oscillation of the
Christoffel/Wronskian profile $K_n(x,x)/r(x)^2$.  For a
Chebyshev envelope that profile varies by $\sim 6\%$ across
nine mask nodes (arcsine at $x=0.6$); the natural OP
oscillation is period $\sim 2$ mesh steps (interior versus
crossing), which produces mixed run lengths $\{1,2,3\}$, not a
one-sided slow-then-fast block.  A Chebyshev $\theta$-compress
at $+x$ ($19$ tails $(1,1,1)$ among $108$ warp-windows)
produces $0$ $V_2$-violators: gradual compression yields mixed
$\mathrm{prev4}$, $V_2$-regular of the same class as $\chi_3$
index $16$.

\section{Adversaries against the source}\label{sec:adv}

\begin{lemma}[not a theorem of Jacobi-positive chains]\label{lem:not-abstract}
There exists a Jacobi-positive coefficient perturbation of a
real window that realises the $v_2$-violator.  Concretely: on
$\chi_3$ at $w=9$ ($N=184$), multiplying the last twelve
$\gamma_k$ by $8$ produces $v_2$-runs ending
$(3,1,3,3,1,1,1)$, which fails $V_2$.  The same window with
scale $6$, or with last $4$, $8$, or $20$ coefficients scaled
by any tested factor in $\{1.2,\ldots,12\}$, does not violate.
\end{lemma}

The source last-eight consecutive $\lvert\log(\gamma_{k+1}/\gamma_k)\rvert$
is at most $0.25$ (ratio $\le 1.28$), and
$\max_k\lvert\gamma_k-\tfrac14\rvert\le 0.077$ on the pins
(last three $\le 0.041$).  The adversarial scale $8$ sits far
outside those bounds.  $V_2$ therefore uses the source, as
r365 demanded; it is not abstract product combinatorics
(already r365) and not abstract Jacobi-positivity.

\begin{lemma}[weight-cluster does not fire]\label{lem:weight-adv}
On FRAME-A $w=9$, depleting $n_{\mathrm{gap}}=6$ arithmetic
atoms just before the $+x$ mask and boosting
$n_{\mathrm{spk}}=3$ atoms at the mask, at boosts $\{8,64\}$,
never produces a $v_2$ or colouring violator of $V_2$.  Both
perturbations break Jacobi positivity ($\gamma_k<0$ for some
$k\le N-2$).
\end{lemma}

The natural source-side recipe for a density step (a mass dip
then a cluster at the mask) is incompatible with $L^*$
positivity on this window, and even after positivity is lost
the occupancy is not the violator (the tail stays a Chebyshev
stub of type $(\ldots,6,2,1)$).  Arithmetic weight ratios at
the mask are already wild (consecutive ratios up to $25$--$70$
on the pins); the OP zeros are more regular than the atoms.

\section{The smaller lemma \texorpdfstring{$V_3'$}{V3'}}\label{sec:v3}

\begin{definition}[$V_3'$]\label{def:v3}
Let $L=15$ and let $(\Delta\theta_i)_{i=0}^{L-2}$ be the
spatial steps of the degree-accumulated Jacobi--Pr\"ufer phase
of $v_2$ on the last $L$ bulk-border nodes at the $+x$
$20\%$ mask, each reduced to $(-\pi,\pi]$ and oriented so that
$\sum\Delta\theta_i>0$.  Write $\mathcal{A}_L\subset\R^{L-1}$
for the region of sequences such that for every incoming
remainder $\eta\in[0,\pi)$, the occupancy of the walk
$\theta_0=\eta$, $\theta_{i+1}=\theta_i+\Delta\theta_i$ is
$V_2$-regular: if it ends $(1,1,1)$, then among the four
preceding runs at least two are $\le 2$.
Condition $V_3'$ is: the source $14$-tuple lies in
$\mathcal{A}_{15}$.
\end{definition}

\begin{prop}[reduction]\label{prop:v3}
$V_3'$ is strictly more elementary than $V_2$ in the sense of
the hydra ward of~\cite{v2}~\S5:
\begin{enumerate}
\item \emph{Finite.}  A point of $\R^{14}$ and a compact scan
  of $\eta\in[0,\pi)$.  $V_2$ is an implication over a
  variable-length run history of an $O(N)$-cell product
  colouring.
\item \emph{Closer to the construction.}  The coordinates are
  the Wronskian steps of the Jacobi chain on the mask mesh
  (Lemma~\ref{lem:wronskian}); no pairing, no $w$, no
  MED-CAP.  The $15$ nodes are the tent/mask geometry.
\item \emph{Fewer quantifiers.}  $\forall\eta$ on a compact
  interval, of a $14$-tuple, versus four preceding runs of
  unbounded length in the product $\mathrm{sign}(v_2 w x)$.
\end{enumerate}
$V_3'$ implies $V_2$ for $v_2$-runs: the actual remainder is
one $\eta$.  At the $+x$ mask, Lemma~\ref{lem:xconst} reduces
the colouring to $\mathrm{sign}(w)\,\mathrm{sign}(v_2)$.  On
MAIN ($w$ constant on the bulk) this is $v_2$-only.  On a
$w$-burst, Chebyshev times Nyquist never violates
(r365 Lemma~6, re-gated).  On the pin windows every
$(1,1,1)$ produced by an $\mathcal{A}_{15}$-profile under
remainder variation is $V_2$-regular; on $\chi_3$ index $16$
the actual occupancy is $\mathrm{prev4}=(3,2,1,3)$ and the
$\eta$-scan has two regular $\mathrm{prev4}$ types.  The chain is
\[
T_2'
 \;\Longleftarrow\;
 g_i\ge\tfrac38
 \;\Longleftarrow\;
 \mathrm{MED\text{-}CAP}
 \;\Longleftarrow\;
 X_n
 \;\Longleftarrow\;
 V_2
 \;\Longleftarrow\;
 V_3'.
\]
\end{prop}

$V_3'$ is logically stronger than $V_2$ for $v_2$-runs (all
remainders, not just the actual one).  That is the content of
the reduction: the occupancy is an alignment, the obstruction
is the speed profile.  Lemma~\ref{lem:contrast} identifies
$\R^{14}\setminus\mathcal{A}_{15}$ with slow-then-fast
contrast blocks.  The remaining lemma is that the source
Wronskian profile at the mask has no such block.

\begin{prop}[Path A on the pins]\label{prop:pins}
On the six pin windows (four $w=9$ worlds, $\chi_4$ index $53$,
$\chi_3$ index $16$), and on four extra FRAME-A rungs
($kz\in\{15,21,25,33\}$):
\begin{enumerate}
\item the last-$15$ Pr\"ufer-step tuple lies in
  $\mathcal{A}_{15}$ ($0$ remainder-violators among $4001$
  remainders at $L\in\{9,13,17\}$);
\item every remainder that produces a tail $(1,1,1)$ does so
  with mixed $\mathrm{prev4}=(3,2,1,3)$ (the $\chi_3$
  index-$16$ type), never with two preceding runs $\ge 3$;
\item colouring $V_2$ holds (regression of~\cite{v2});
\item $\chi_3$ index $16$ is the $v_2$-triple, even run count,
  $n_0=2$; $\chi_4$ index $53$ is the $w$-driven triple,
  odd, saturating prefix $(1,2,6,5,4,4)$.
\end{enumerate}
This is a pin census of $V_3'$, not a cofinal law.
\end{prop}

The r365 hydra candidates (mesh-only $v_2$; $2$-regular
$v_2$; $w$-only; Freud / $\gamma_k\to 1/4$) fail the elementary
test or are false.  $V_3'$ does not: it is a finite region in
the Wronskian coordinates of a $15$-node mask window.  It does
not reduce toward L$^*$ (no coefficient-asymptotic is assumed;
Lemma~\ref{lem:not-abstract} shows that a large coherent
$\gamma$-scale, far outside the source, \emph{does} leave
$\mathcal{A}_{15}$).

\section{What remains}\label{sec:rest}

To prove $V_3'$ cofinally: show that the source Wronskian
profile $K_n(x,x)/r(x)^2$ on the last $15$ bulk nodes has no
slow-then-fast block of the Lemma~\ref{lem:contrast} type.
The envelope of a near-Chebyshev Jacobi chain varies too slowly
at $x\approx 0.6$ (relative change $O(1/N)$ across the window);
the OP oscillation is two-step, not a one-sided burst.  The
coefficient bounds of the source ($\lvert\gamma-1/4\rvert\le 0.077$
on the pins, consecutive ratio $\le 1.28$) exclude the only
explicit coefficient adversary that left $\mathcal{A}_{15}$
(Lemma~\ref{lem:not-abstract}).  Making that comparison a
theorem --- a quantitative Sturm bound on $K_n(\,\cdot\,,\,\cdot\,)$
at the mask, from the tent/mesh and the von~Mangoldt weights
--- is the remaining work.  It is a statement about fifteen
nodes and the Wronskian, not about $V_2$.

T1 / $C_K$ unchanged.  The cubic-mass bound of~\cite{medcap}
still carries $C_K$.  No RH claim.

\section{Machine checks}\label{sec:machine}

Every numbered lemma is re-proved in
\texttt{verify\_v2\_lemma.py} (same directory).  Standalone:
Lagrange identity over $\Q$; contrast dichotomy; Chebyshev
mask and $\theta$-compress; the named combinatorial violator.
Construction: Wronskian residual, Path~A on the pins, $\gamma$
bounds, the $8\times$ last-$12$ adversary, $+x$ sign, the
r372 dictionary on the pins, colouring $V_2$ regression, the
weight-cluster pair at boosts $8$ and $64$.
Exit line: \texttt{V2 LEMMA V3 VERIFIED}.

\section*{Claim boundary}

Finite identities on an explicit Wronskian, a finite contrast
dichotomy in $\R^{24}$, a named coefficient adversary, a
weight-cluster negative, and a strictly smaller finite profile
lemma $V_3'$ to which $V_2$ reduces.  No statement about zeros
of $\zeta(s)$, in either direction.  No RH claim.

\begin{thebibliography}{9}
\bibitem{medcap}
S.~Hamann,
\textit{A median cap for gaps of a tiled integer packing},
standalone note, August 2026,
\texttt{rh/problem/medcap\_lemma.tex}.
\bibitem{xn}
S.~Hamann,
\textit{The length invariant $X_n$ of a folded sign field},
standalone note, August 2026,
\texttt{rh/problem/xn\_invariant.tex}.
\bibitem{v2}
S.~Hamann,
\textit{Run-length regularity $V_2$ at the $x$-mask},
standalone note, August 2026,
\texttt{rh/problem/v2\_regularity.tex}.
\end{thebibliography}

\end{document}
